\section{可微加权三角化方法}

\subsection{研究目的与动机}

在双目人体三维姿态估计任务中，我们的目标是利用多个视角下的二维关键点观测，恢复对应的人体三维关节坐标。
在相机参数已知的条件下，传统多视图几何中的三角化（Triangulation）可以通过解析几何方法求解三维点位置。

然而，在深度学习框架下直接使用传统三角化存在两个关键问题：

\begin{enumerate}
    \item \textbf{二维关键点检测存在噪声。} 由于遮挡、模糊或复杂姿态，某些视图中的二维关键点可能存在较大误差。
    \item \textbf{不可端到端训练。} 若通过 $\arg\max$ 从热力图中提取二维坐标，则该操作不可导，三维监督信号无法反向传播至热力图网络。
\end{enumerate}

为解决上述问题，我们提出一种\textbf{可微加权三角化（Differentiable Weighted Triangulation）}方法。该方法具有以下特性：

\begin{itemize}
    \item 使用 soft-argmax 从热力图中提取二维坐标，保证可微性；
    \item 为每个视图分配可学习置信度权重，增强鲁棒性；
    \item 将三角化问题转化为加权最小二乘问题；
    \item 通过特征值分解求解三维点，保证端到端可训练。
\end{itemize}

\subsection{符号定义}

设：

\begin{itemize}
    \item $C$ 表示视图数量（双目时 $C=2$）；
    \item $j$ 表示关节索引；
    \item $\mathbf{X}_j = [X_j, Y_j, Z_j]^T \in \mathbb{R}^3$ 为三维关节坐标；
    \item $\tilde{\mathbf{X}}_j = [X_j, Y_j, Z_j, 1]^T \in \mathbb{R}^4$ 为其齐次形式；
    \item $\mathbf{x}_{c,j} = [u_{c,j}, v_{c,j}]^T$ 为第 $c$ 个视图下的二维关键点；
    \item $\tilde{\mathbf{x}}_{c,j} = [u_{c,j}, v_{c,j}, 1]^T$ 为其齐次形式；
    \item $\mathbf{P}_c \in \mathbb{R}^{3\times4}$ 为第 $c$ 个相机的投影矩阵；
    \item $w_{c,j} > 0$ 为第 $c$ 个视图对关节 $j$ 的置信度权重。
\end{itemize}

设投影矩阵 $\mathbf{P}_c$ 的三行分别为：

\[
\mathbf{P}_c =
\begin{bmatrix}
\mathbf{p}_{1,c}^T \\
\mathbf{p}_{2,c}^T \\
\mathbf{p}_{3,c}^T
\end{bmatrix}
\]

\subsection{投影约束与线性化}

根据透视投影模型，在齐次坐标下有：

\[
\tilde{\mathbf{x}}_{c,j} \sim \mathbf{P}_c \tilde{\mathbf{X}}_j
\]

其中“$\sim$”表示两向量仅差一个尺度因子。
为了消除尺度因子，引入叉乘约束：

\[
\tilde{\mathbf{x}}_{c,j} \times (\mathbf{P}_c \tilde{\mathbf{X}}_j) = 0
\]

展开后可得两条线性独立方程：

\[
(u_{c,j} \mathbf{p}_{3,c}^T - \mathbf{p}_{1,c}^T)\tilde{\mathbf{X}}_j = 0
\]

\[
(v_{c,j} \mathbf{p}_{3,c}^T - \mathbf{p}_{2,c}^T)\tilde{\mathbf{X}}_j = 0
\]

因此，第 $c$ 个视图对关节 $j$ 提供两个线性约束，可写为：

\[
\mathbf{A}_{c,j}\tilde{\mathbf{X}}_j = 0
\]

其中

\[
\mathbf{A}_{c,j} =
\begin{bmatrix}
u_{c,j} \mathbf{p}_{3,c}^T - \mathbf{p}_{1,c}^T \\
v_{c,j} \mathbf{p}_{3,c}^T - \mathbf{p}_{2,c}^T
\end{bmatrix}
\]

\subsection{多视图线性系统}

将所有视图堆叠，可得整体线性系统：

\[
\mathbf{A}_j \tilde{\mathbf{X}}_j = 0,
\quad
\mathbf{A}_j =
\begin{bmatrix}
\mathbf{A}_{1,j} \\
\mathbf{A}_{2,j} \\
\vdots \\
\mathbf{A}_{C,j}
\end{bmatrix}
\in \mathbb{R}^{2C \times 4}
\]

由于存在噪声，一般不存在精确解，因此转化为最小二乘问题：

\[
\min_{\|\tilde{\mathbf{X}}_j\|=1}
\| \mathbf{A}_j \tilde{\mathbf{X}}_j \|^2
\]

\subsection{加权最小二乘形式}

为提高对噪声的鲁棒性，引入置信度权重 $w_{c,j}$。
构造对角权重矩阵：

\[
\mathbf{W}_j =
\mathrm{diag}(
\sqrt{w_{1,j}}, \sqrt{w_{1,j}},
\sqrt{w_{2,j}}, \sqrt{w_{2,j}},
\ldots,
\sqrt{w_{C,j}}, \sqrt{w_{C,j}}
)
\]

加权线性系统为：

\[
\mathbf{B}_j = \mathbf{W}_j \mathbf{A}_j
\]

优化问题变为：

\[
\min_{\|\tilde{\mathbf{X}}_j\|=1}
\| \mathbf{B}_j \tilde{\mathbf{X}}_j \|^2
\]

\subsection{特征值求解}

展开目标函数：

\[
\| \mathbf{B}_j \tilde{\mathbf{X}}_j \|^2
=
\tilde{\mathbf{X}}_j^T
(\mathbf{B}_j^T \mathbf{B}_j)
\tilde{\mathbf{X}}_j
\]

记：

\[
\mathbf{M}_j = \mathbf{B}_j^T \mathbf{B}_j
\]

则问题转化为：

\[
\min_{\|\tilde{\mathbf{X}}_j\|=1}
\tilde{\mathbf{X}}_j^T \mathbf{M}_j \tilde{\mathbf{X}}_j
\]

根据 Rayleigh 商定理，最优解为 $\mathbf{M}_j$ 最小特征值对应的特征向量：

\[
\mathbf{M}_j \tilde{\mathbf{X}}_j = \lambda \tilde{\mathbf{X}}_j
\]

取最小特征值对应的特征向量作为 $\tilde{\mathbf{X}}_j$。

\subsection{反齐次化}

得到齐次解：

\[
\tilde{\mathbf{X}}_j =
\begin{bmatrix}
X_h \\
Y_h \\
Z_h \\
W_h
\end{bmatrix}
\]

则欧式三维坐标为：

\[
\mathbf{X}_j =
\frac{1}{W_h}
\begin{bmatrix}
X_h \\
Y_h \\
Z_h
\end{bmatrix}
\]

\subsection{与热力图的可微连接}

二维关键点由热力图通过 soft-argmax 获得：

\[
u_{c,j} = \sum_{x,y} x \cdot
\mathrm{softmax}(H_{c,j}(x,y))
\]

\[
v_{c,j} = \sum_{x,y} y \cdot
\mathrm{softmax}(H_{c,j}(x,y))
\]

由于 softmax 与加权求和均为可微操作，故
$\mathbf{X}_j$ 关于热力图 $H_{c,j}$ 可导，
从而三维监督损失能够端到端反向传播。

\subsection{总结}

本文将传统三角化问题重写为加权最小二乘优化问题，
并通过特征值分解求解三维点。
通过 soft-argmax 与可学习置信度权重的引入，
实现了三维几何约束与深度学习网络的无缝结合，
构建了一个完全可微、鲁棒且端到端可训练的多视图三维姿态估计框架。